\section{Síntesis y ampliación}

\subsection{El cambio de paradigma en la investigación}

Yang Zhilin (杨植麟) señala que el paradigma de investigación en IA actual es radicalmente distinto al de hace 10 años. Antes se privilegiaba «publicar una nueva idea», pero sin el respaldo de experimentos rigurosos era difícil llegar a conclusiones confiables. Ahora, con la scaling ladder y un rico conjunto de benchmarks, es posible validar las ideas en distintas escalas, lo que permite hacer mejoras seguras y sólidas sobre técnicas «antiguas» (como los optimizadores, el mecanismo de atención y las conexiones residuales).

\subsection{El efecto multiplicativo de las tres dimensiones}

\begin{importantbox}{De la suma a la multiplicación}
Las tres dimensiones de scaling no se suman de manera simple, sino que mantienen una relación multiplicativa de refuerzo mutuo:
\begin{itemize}
    \item Adam (2014) $\rightarrow$ Muon-Clip: mejor token efficiency
    \item Full Attention (2017) $\rightarrow$ Kimi Linear: mejor contexto largo
    \item Residual Connection $\rightarrow$ Attention Residue: scaling de profundidad más eficiente
\end{itemize}
Al multiplicar estas mejoras, se obtiene una ganancia global muy superior a la suma de las mejoras individuales.
\end{importantbox}

\subsection{El futuro de la comunidad de código abierto}

Yang Zhilin se muestra optimista sobre el futuro de la comunidad de código abierto y considera que seguirán surgiendo mejoras revolucionarias en arquitectura y optimización. Agent Swarms no es el punto final: seguirán apareciendo nuevas dimensiones de scaling, lo que impulsará una competencia continua entre los modelos de código abierto y los de código cerrado.

\subsection{Lecturas adicionales}

\begin{itemize}
    \item Kaplan et al., «Scaling Laws for Neural Language Models» (2020) --- scaling law clásica
    \item Artículo de Muon Optimizer --- trabajo del equipo de Kimi sobre un optimizador de segundo orden
    \item Reporte técnico de Kimi Linear --- derivación detallada de las fórmulas de Delta Attention
    \item Reporte técnico de Attention Residue --- experimentos completos sobre la atención en la dimensión de profundidad
    \item He et al., «Deep Residual Learning for Image Recognition» (2015) --- artículo original de ResNet
\end{itemize}

%% --- Fin del contenido --- %%

\end{document}
